\documentclass[11pt,oneside]{book}
\usepackage[english]{babel} 
\usepackage{microtype}
\usepackage[utf8]{inputenc}
\usepackage{subfigure}
\usepackage{epsfig}
\usepackage{amsmath, amssymb}
\usepackage[linesnumbered,lined,commentsnumbered,italiano]{algorithm2e}
\usepackage{fancybox}
\usepackage{listings}
\usepackage{url}
\usepackage{graphicx}
\usepackage{todonotes}
\usepackage{moreverb}
\usepackage[labelfont=it,textfont={it}]{caption}
\usepackage[final,bookmarks,breaklinks,colorlinks,allcolors=black]{hyperref}
\usepackage{color}
\usepackage[block=ragged,backend=biber,natbib=true]{biblatex} % Use the bibtex backend with the authoryear citation style (which resembles APA)
\addbibresource{sections/bibliography.bib} % The filename of the bibliography
\usepackage[autostyle=true]{csquotes} % Required to generate language-dependent quotes in the bibliography


\definecolor{mygreen}{rgb}{0,0.6,0}
\definecolor{mygray}{rgb}{0.5,0.5,0.5}
\definecolor{mymauve}{rgb}{0.58,0,0.82}
\lstset{ %
	backgroundcolor=\color{white},   % choose the background color; you must add \usepackage{color} or \usepackage{xcolor}; should come as last argument
	basicstyle=\footnotesize,        % the size of the fonts that are used for the code
	breakatwhitespace=false,         % sets if automatic breaks should only happen at whitespace
	breaklines=true,                 % sets automatic line breaking
	captionpos=b,                    % sets the caption-position to bottom
	commentstyle=\color{mygreen},    % comment style
	deletekeywords={...},            % if you want to delete keywords from the given language
	escapeinside={\%*}{*)},          % if you want to add LaTeX within your code
	extendedchars=true,              % lets you use non-ASCII characters; for 8-bits encodings only, does not work with UTF-8
	frame=single,
	keepspaces=true,                 % keeps spaces in text, useful for keeping indentation of code (possibly needs columns=flexible)
	keywordstyle=\color{blue},       % keyword style
	morekeywords={*,...},            % if you want to add more keywords to the set
	numbers=left,                    % where to put the line-numbers; possible values are (none, left, right)
	numberstyle=\tiny\color{mygray}, % the style that is used for the line-numbers
	rulecolor=\color{black},         % if not set, the frame-color may be changed on line-breaks within not-black text (e.g. comments (green here))
	showspaces=false,                % show spaces everywhere adding particular underscores; it overrides 'showstringspaces'
	showstringspaces=false,          % underline spaces within strings only
	showtabs=false,                  % show tabs within strings adding particular underscores
	stringstyle=\color{mymauve},     % string literal style
	tabsize=2,	                   % sets default tabsize to 2 spaces
	title=\lstname                   % show the filename of files included with \lstinputlisting; also try caption instead of title
}


\definecolor{gray}{rgb}{0.4,0.4,0.4}
\definecolor{darkblue}{rgb}{0.0,0.0,0.6}
\definecolor{cyan}{rgb}{0.0,0.6,0.6}


\definecolor{javared}{rgb}{0.6,0,0} % for strings
\definecolor{javagreen}{rgb}{0.25,0.5,0.35} % comments
\definecolor{javapurple}{rgb}{0.5,0,0.35} % keywords
\definecolor{javadocblue}{rgb}{0.25,0.35,0.75} % javadoc


\lstdefinelanguage{FLY}
{
  morestring=[b]",
  morecomment=[l]{//},
  morecomment=[s]{/}{/},
  identifierstyle=\color{black},
  keywordstyle=\color{javapurple},
  morekeywords={func,var,while, if, then, else,for, in, println, fly ,on,thenall,as,then,require,native}% list your attributes here
}


\lstset{language=FLY,
basicstyle=\ttfamily\scriptsize,
keywordstyle=\color{javapurple}\bfseries,
stringstyle=\color{javared},
commentstyle=\color{javagreen},
morecomment=[s][\color{javadocblue}]{/*}{/},
numbers=left,
numberstyle=\tiny\color{black},
stepnumber=1,
numbersep=10pt,
tabsize=4,
showspaces=false,
showstringspaces=false}


\begin{document}
\begin{titlepage}
\begin{center}
\epsfig{file=figures/logo_standard.jpg,width=2.5truecm}\\[0.2truecm]
{\Large Universit\`a degli Studi di Salerno}\\[0.2truecm]
{\large Dipartimento di Informatica}\\
\hrulefill
\vfill
{\large Tesi di Laurea di I livello in }\\[0.2truecm]
{\Large Informatica}\\
\vfill\vfill
{\Huge Template tesi ISISLab}
\vfill\vfill


{\bf Relatore} \hfill {\bf Candidato}\ \ \\
Nome Cognome \hfill Nome Cognome\\
{\bf Correlatore} \hfill {\bf }\ \ \\
Dott. Nome Cognome \hfill \ \ \\

\vfill
\hrulefill 

Academic Year 2021-2022

\end{center}
\end{titlepage}

\pagenumbering{roman}
\chapter*{Abstract}
Lorem ipsum dolor sit amet, consectetur adipiscing elit, sed do eiusmod tempor incididunt ut labore et dolore magna aliqua. Ut enim ad minim veniam, quis nostrud exercitation ullamco laboris nisi ut aliquip ex ea commodo consequat. Duis aute irure dolor in reprehenderit in voluptate velit esse cillum dolore eu fugiat nulla pariatur. Excepteur sint occaecat cupidatat non proident, sunt in culpa qui officia deserunt mollit anim id est laborum.



\tableofcontents
\pagestyle{plain}

%%%%%%%%%%%%%%%%%%%%%%%%%%%%
\chapter{Introduction}
\setcounter{page}{1} 	% devono seguire solo il primo capitolo
\pagenumbering{arabic}	% devono seguire solo il primo capitolo
A computer hardware and software combination created for a particular purpose is an Embedded System (\es). In many cases, \ess operate as part of a bigger system (\eg  agricultural and processing sector equipment, automobiles, medical equipment, airplanes, and so on) and within tight resource constrains (\eg small battery, limited memory and CPU speed, and so on).
The global \ess market is expected to witness notable growth. A recent report evaluated the  global \ess market 89.1 billion dollars in 2021, and this market is projected to reach 163.2 billion dollars by 2031; with a compound annual growth rate of  6.5\%~ \cite{ESSTR2022}. This growth is mostly related to  an increase in the demand for advanced driver-assistance system (in electric and hybrid vehicles) and in the number of \ess-related research and development projects.  


Test-Driven Development (\tdd) is an incremental approach to software development in which a developer repeats a short development cycle made up of three phases: \textit{red}, \textit{green}, and \textit{blue/refactor} \cite{TDDByExample}. During the red phase, the developer writes a unit test for a small chunk of a functionality not yet implemented and watches the test fail. In the green phase, the developer implements the chunk of functionalities as quickly as possible and watches all the unit tests pass. During the refactoring phase, the developer changes the internal structure of the code while paying attention to keep the functionality intact—accordingly, all unit tests should pass. TDD has been conceived to develop “regular” software, and it is claimed to improve software quality as well as developers' productivity \cite{DBLP:reference/se/ErdogmusMJ10}. \ess have all the same challenges of non-embedded systems (\noess), such as poor quality, but add challenges of their own [6]. 
For example, one of the most cited differences between embedded and \noess is that embedded code depends on the hardware. While in principle, there is no difference between a dependency on a hardware device and one on a \noess \cite{TDDEC}, dealing with hardware introduces a whole new set of variables to consider during development. Furthermore, the limited resources on which an \es usually operates may make it extremely difficult to properly test the system in its entirety. 
Over the years, a huge amount of empirical investigations has been conducted to study the claimed effects of TDD on the development of \noess (\eg, \cite{DBLP:journals/software/KaracT18}). So far no investigations have been conducted to assess possible benefits concerned the application of TDD on the development of \ess although some authors, like Greening \cite{TDDEC}, believes that embedded developers can benefit from the application of TDD in the development of \ess.

In this work of thesis, we investigate the following primary research question (RQ):
\\ \ \\
\noindent \textbf{RQ.} Does the use of \tdd affect the development of \ess?	
\\ \ \\
To answer this RQ, we present the results of a controlled (baseline) experiment and its replication to study the application of \tdd on the implementation of \ess. The participants were final year Master Students in Computer Science enrolled to an \es course at the University of Salerno. The goal of these experiments is to increase our body of knowledge on the benefit (if any) related to the application of \tdd in the context of the development of \ess. To that end, we compare \tdd with respect to a more traditional, test-last, development practice, where test cases are written after the production code. From here onwards, we refer to this traditional way of coding as \notdd. Whichever the used approach (\tdd and \notdd), the participants in the implementation of an \es where asked to use a mock to model and confirm the interactions between a device driver and the hardware. The mocked implementation would intercept commands to and from the device simulating a given usage scenario. In the replication experiment, the participants had to implement a more complex \es and had also to replace the mocks with real hardware components (a number of sensors and actuators) before deploying the \es on the actual hardware platform they had mocked up to that point. The logic of the \es was deployed on Raspberry Pi model 3 and the developed test cases were executed in the real environment. Variations in the replicated experiments (task and the experimental procedure) were introduced to validate the results of the baseline experiment and to generalize these results to a more real setting. The results obtained by an aggregated analysis of the data from both the experiments suggest that there are no significant differences between \tdd and \notdd on: productivity, number of written tests, a..., respectively. On the other hand, we observed a significant difference on: ... When analyzing experiments individually, we observed that ...

As for the following thesis structure, it is made up of five chapters:
The first two act as the foundational knowledge concepts that provide a general overview on the main topics surrounding \tdd and \es.
More in detail, chapter one contains an overview on the software testing process, analyzing the main approaches, with a focus on \tdd. Chapter two will provide information on the general \ess concepts, enabling technologies and implementation challenges, before discussing the techniques for testing such systems.
In chapter three, we conduct a review of the literature by examining the most relevant previous empirical studies on the application of \tdd and the current testing methodologies for \es.
Chapter four will contain the detail explanation of our approach for the definition of the two experimental studies, the analysis of the results, and our answers to the research questions.
Finally, chapter five will provide the conclusions to the thesis, along with a discussion on the possible ramification the research could manifest when moving forward in the analysis of \tdd for \es development.
 
%%%%%%%%%%%%%%%%%%%%%%%%%%%%
\chapter{Literature}
\section{Studies on Test Driven Development}
The Empirical Software Engineering community has taken interest into the investigation of the effects of TDD on several efforts, including testing efforts, external software quality and developers productivity \cite{DBLP:conf/esem/FucciS0SSUTJO16} \cite{DBLP:journals/tse/ErdogmusMT05} \cite{DBLP:journals/infsof/Madeyski10}.

However, the empirical evidence so far has been mixed regarding the effects of TDD: \dots



\section{TDD for Embedded Systems}
Often, quality in embedded software is generally tied to platform-specific testing tools geared towards debugging \cite{TDDEmbeddedSoftware}.
\dots

Applying TDD practices to ESs could potentially result in a series of benefits:
\begin{itemize}
    \item \dots
\end{itemize}


\begin{itemize}
    \item Reduce risk by verifying production code, independent of hardware, before hardware is ready or when hardware is expensive and scarce.
    \item Reduce the number of long target compile, link, and upload cycles that are executed by removing bugs on the development system.
    \item Reduce debug time on the target hardware where problems are more difficult to find and fix.
    \item Isolate hardware/software interaction issues by modeling hardware interactions in the tests.
    \item Improve software design through the decoupling of modules from each other and the hardware. Testable code is by necessity, modular.
\end{itemize}


\noindent In \cite{TDDEC}, the author proposes the "Embedded TDD Cycle", as a pipeline made of the following steps:
\begin{enumerate}
    \item \textbf{TDD micro-cycle}: this first stage is the one run most frequently, usually every few minutes. During this stage, a bulk of code is written in TDD fashion, and compiled to run on the host development system: doing so gives the developer fast feedback, not encumbered by the constraints of hardware reliability and/or availability, since there are no target compilers or lengthy upload processes. Furthermore, the development system should be a proven and stable execution environment, and usually has a richer debugging environment compared to the target platform. 
    Running the code on the development system, when it is eventually going to run in a foreign environment can be risky, so it's best to confront that risk regularly.
    \item \textbf{Compiler Compatibility Check}: periodically compile for the target environment, using the cross-compiler expected to be used for production compilations; this stage can be seen as an early warning system for any compiler incompatibilities, since it warns the developer of any porting issue, such as unavailable header files, incompatible language support, and missing language features. As a result, the written code only uses facilities available in both development environments.
    A potential issue at this stage is that in early ES development, the tool chain may not yet be decided, and this compatibility check cannot be performed: in this case, developers should take their best guess on the tool chain and compile against that compiler.
    Finally, this stage should not run with every code change; instead, a target cross-compile should take place whenever a new language feature is used, a new header file is included or a new library call is performed.
    \item \textbf{Run unit tests in an evaluation board}: compiled code could potentially run differently in the host development system and the target embedded processor. In order to mitigate this risk, developers can run the unit tests on an evaluation board; with this, any behavior differences between environments would emerge, ad since runtime libraries are usually prone to bugs \cite{TDDEC}, the risk is real. If it's late in the development cycle, and a reliable target hardware is available, this stage may appear unnecessary.
    \item \textbf{Run unit tests in the target hardware}: the objective here is again to run the test suite, however this time doing so while exercising the real hardware. One additional aspect to this stage is that developers could also run target hardware-specific tests. These tests allow developers to characterize or learn how the target hardware behaves. An additional challenge in this stage is limited memory in the target. The entire unit test suite may not fit into the target. In that case, the tests can be reorganized into separate test suites, where each suite fits in memory. This, however, does result in more complicated build automation process.
    \item \textbf{Run acceptance tests in the target hardware}: Finally, in order to make sure that the product features work, automated and manual acceptance tests are run in the target environment. Here developers have to make sure that any of the hardware-dependent code that can't be fully tested automatically is tested manually.
\end{enumerate}

%%%%%%%%%%%%%%%%%%%%%%%%%%%%
\chapter{Conclusions}
Lorem ipsum dolor sit amet, consectetur adipiscing elit, sed do eiusmod tempor incididunt ut labore et dolore magna aliqua. Ut enim ad minim veniam, quis nostrud exercitation ullamco laboris nisi ut aliquip ex ea commodo consequat. Duis aute irure dolor in reprehenderit in voluptate velit esse cillum dolore eu fugiat nulla pariatur. Excepteur sint occaecat cupidatat non proident, sunt in culpa qui officia deserunt mollit anim id est laborum.

%%%%%%%%%%%%%%%%%%%%%%%%%%%%
\nocite{*}
\printbibliography[title={Bibliography}] 

%%%%%%%%%%%%%%%%%%%%%%%%%%%%
\end{document}
